% cv-entry-inline-order.tex — extraction fixture for `entrymeta=inline'.
%
% Purpose (issue #230)
%   The CV half of the demonstration. resume-entry-inline-order.tex carries the
%   full explanation of the mechanism; this file exists because the option is
%   declared on both record classes and the acceptance criterion is stated for
%   both, and because the CV reaches the same component through its own
%   \CDossierEntry with its own surrounding spacing.
%
%   This is cv-entry-dates-order.tex with `entrymeta=inline' and
%   `\CDossierRecordListEdgeAboveSkip' at 0.125, half its committed floor. The
%   column baseline extracts each entry's dates and location as a separate block
%   after the heading; here they extract on the heading lines, joined by `|', in
%   source order.
%
%   The 0.125 override is the assertion. Do not raise it to make a failure go
%   away — see resume-entry-inline-order.tex.
%
%   The committed baseline (cv-entry-inline-order.expected.txt) is the
%   assertion; regenerate it only for an intended, reviewed output change.
%
% Compile: lualatex -halt-on-error -interaction=nonstopmode cv-entry-inline-order.tex
% (run from the repository root, or with the root on TEXINPUTS as run.sh does)
\documentclass[entrymeta=inline]{careerdossier-cv}

\ExplSyntaxOn
\skip_set:Nn \CDossierRecordListEdgeAboveSkip
  { \fp_to_dim:n { 0.125 * \dim_to_fp:n { \CDossierBodyLeading } } }
\ExplSyntaxOff

\CDossierSetup{
  name     = {Ada Lovelace},
  headline = {Research Engineer},
  email    = {ada@example.edu},
  phone    = {+1 416 555 0142},
  website  = {example.edu/ada},
}
\begin{document}
\MakeCDossierHeader
\CDossierSection{Research Experience}
\begin{CDossierEntry}[
  title        = {Postdoctoral Researcher},
  organization = {Example University},
  location     = {Toronto, ON},
  dates        = {2024--Present}
]
  \begin{CDossierItemize}
    \item First research contribution in source order.
    \item Second research contribution with a reproducible analysis pipeline.
  \end{CDossierItemize}
\end{CDossierEntry}
\begin{CDossierEntry}[
  title        = {Graduate Researcher},
  organization = {Example Institute},
  location     = {Ottawa, ON},
  dates        = {2020--2024}
]
  \begin{CDossierItemize}
    \item Later contribution that follows the first entry on the same page.
  \end{CDossierItemize}
\end{CDossierEntry}
\end{document}
